\documentclass[11pt]{article}

\usepackage[T1]{fontenc}
\usepackage[utf8]{inputenc}
\usepackage{lmodern}
\usepackage[margin=1in]{geometry}
\usepackage{hyperref}
\usepackage{xcolor}
\usepackage{listings}
\usepackage{enumitem}
\usepackage{longtable}

\hypersetup{
  colorlinks=true,
  linkcolor=blue!60!black,
  urlcolor=blue!60!black,
  pdftitle={PlussQit User Guide},
  pdfauthor={PlussMa Contributors}
}

\lstdefinestyle{plusscode}{
  basicstyle=\ttfamily\small,
  keywordstyle=\color{blue!70!black},
  stringstyle=\color{green!40!black},
  commentstyle=\color{gray!80!black},
  numbers=left,
  numberstyle=\tiny\color{gray},
  stepnumber=1,
  numbersep=8pt,
  frame=single,
  breaklines=true,
  showstringspaces=false,
  tabsize=2
}

\lstset{style=plusscode}

\title{PlussQit User Guide\\\large Practical Series/Note/Timeline/TimeIndex Documentation}
\author{PlussMa Contributors}
\date{\today}

\begin{document}
\maketitle
\tableofcontents
\newpage

\section{Purpose}
This guide is focused on \textbf{how to use} PlussQit in real code.
It is organized in this order:
\begin{enumerate}[leftmargin=1.5em]
  \item \texttt{Series} (CodeM)
  \item \texttt{Note} (CodeN)
  \item \texttt{Timeline + Duration} (CodeH)
  \item \texttt{TimeIndex} (CodeH)
\end{enumerate}

For exhaustive API inventories, see Section~\ref{sec:api-index-files}.

\section{Build and Run Demos}
\subsection{Build with examples enabled}
\begin{lstlisting}[language=bash]
cmake --preset core-dev
cmake --build --preset core-dev
\end{lstlisting}

\subsection{Most useful demo binaries}
\begin{itemize}[leftmargin=1.5em]
  \item \texttt{build/PlussQit/plussqit\_codem\_series\_usage\_demo}
  \item \texttt{build/PlussQit/plussqit\_codem\_series\_methods\_demo}
  \item \texttt{build/PlussQit/plussqit\_codem\_stats\_demo}
  \item \texttt{build/PlussQit/plussqit\_codem\_signal\_demo}
  \item \texttt{build/PlussQit/plussqit\_coden\_note\_factory\_demo}
  \item \texttt{build/PlussQit/plussqit\_coden\_note\_method\_families\_demo}
  \item \texttt{build/PlussQit/plussqit\_codeh\_timeline\_non\_generated\_and\_factory\_methods\_demo}
  \item \texttt{build/PlussQit/plussqit\_codeh\_timeline\_dataframe\_timeindex\_methods\_demo}
\end{itemize}

\section{Series (CodeM)}
\subsection{Include and Type Aliases}
\begin{lstlisting}[language=C++]
#include <pluss/pluss.h>
// narrower:
#include <pluss/codem/codem.h>
#include <pluss/codem/series_factory.h>

using pluss::Series;
using pluss::ScalarType;
using pluss::Platform;
\end{lstlisting}

\subsection{Construction Cookbook}
\subsubsection{1) Scalars}
\begin{lstlisting}[language=C++]
auto a = pluss::value(3.14);                  // scalar float
auto b = pluss::value(42);                    // scalar int
auto c = pluss::value(true);                  // scalar bool
auto d = pluss::value(3.14, ScalarType::Float64);
\end{lstlisting}

\subsubsection{2) 1D/2D/N-D from literals}
\begin{lstlisting}[language=C++]
auto v = pluss::value({1.0, 2.0, 3.0});
auto m = pluss::value({{1.0, 2.0}, {3.0, 4.0}});
auto t = pluss::value({{{1,2},{3,4}}, {{5,6},{7,8}}});
\end{lstlisting}

\subsubsection{3) From lvalue host buffers}
\begin{lstlisting}[language=C++]
std::vector<double> flat{1,2,3,4,5,6};
auto x = pluss::make_series(flat, {2,3});
\end{lstlisting}

\subsubsection{4) Shape factories}
\begin{lstlisting}[language=C++]
auto e = pluss::empty({2,3}, ScalarType::Float32);
auto z = pluss::zeros({2,3});
auto o = pluss::ones({2,3});
auto f = pluss::full({2,3}, 7.5);
\end{lstlisting}

\subsubsection{5) Range and random factories}
\begin{lstlisting}[language=C++]
auto a = pluss::arange(0, 10);
auto b = pluss::linspace(0.0, 1.0, 11);
auto c = pluss::randn({3,4});
auto d = pluss::randint(0, 100, {2,3});
auto e = pluss::normal(0.0, 1.0, {5});
\end{lstlisting}

\subsubsection{6) Like-factories}
\begin{lstlisting}[language=C++]
auto base = pluss::zeros({2,3}, ScalarType::Float64);
auto a = pluss::zeros_like(base);
auto b = pluss::rand_like(base, ScalarType::Float64, 123);
auto c = pluss::eye_like(base);
\end{lstlisting}

\subsubsection{7) Complex factories}
\begin{lstlisting}[language=C++]
auto re = pluss::linspace(0.0, 1.0, 4);
auto im = pluss::linspace(1.0, 2.0, 4);
auto z1 = pluss::complex(re, im, ScalarType::Complex64);
auto z2 = pluss::polar(re.abs(), im, ScalarType::Complex64);
\end{lstlisting}

\subsection{Core Usage Patterns}
\subsubsection{Inspect and move data}
\begin{lstlisting}[language=C++]
auto x = pluss::value({1.0,2.0,3.0});
std::cout << x.dtype() << " " << x.sizes() << "\n";

auto cpu_f32 = x.to(pluss::Platform::CPU, pluss::ScalarType::Float32);
auto host = cpu_f32.to_host<float>();
\end{lstlisting}

\subsubsection{Math and reductions}
\begin{lstlisting}[language=C++]
auto x = pluss::normal(0.0, 1.0, {20});
auto y = x.abs().add(1.0).log();
auto mean = y.mean();
auto stdv = y.std(false);
auto s = y.sum();
\end{lstlisting}

\subsubsection{Indexing/shape transforms}
\begin{lstlisting}[language=C++]
auto m = pluss::value({1,2,3,4,5,6}, {2,3});
auto t = m.transpose(0,1);
auto r = m.reshape({3,2});
auto f = m.flatten(0,1);
\end{lstlisting}

\subsubsection{Masking and gather/scatter}
\begin{lstlisting}[language=C++]
auto src = pluss::value({10,20,30,40,50,60}, {2,3});
auto idx = pluss::value({2,1,0,0,1,2}, {2,3});
auto g = src.gather(1, idx);

auto base = pluss::zeros({2,3}, ScalarType::Int64);
auto val  = pluss::value({1,2,3,4,5,6}, {2,3});
auto out  = base.scatter_add(1, idx, val);
\end{lstlisting}

\subsubsection{Signal/statistics utilities}
\begin{lstlisting}[language=C++]
auto x = pluss::value(std::vector<double>{-3,-2,-1,0,1,2,3}, {7});
auto xdct = x.dct();
auto xrec = xdct.idct();
auto h = x.hilbert();
auto hist = x.histcounts(8, -3.0, 3.0);
auto skew = x.skewness();
auto kurt = x.kurtosis();
\end{lstlisting}

\subsection{Series Demo Map}
\begin{longtable}{@{}p{0.36\linewidth}p{0.56\linewidth}@{}}
\textbf{Demo} & \textbf{Best for} \\
\hline
\texttt{series\_usage\_demo.cpp} & Core construction, device/dtype, printing, host extraction \\
\texttt{series\_methods\_demo.cpp} & Broad method families with pass/fail output \\
\texttt{stats\_demo.cpp} & Moments, distributions, moving windows, NaN-aware stats \\
\texttt{signal\_demo.cpp} & DCT/DST/Hilbert/phase/fftshift workflows \\
\end{longtable}

\section{Note (CodeN)}
\subsection{Include and Basic Construction}
\begin{lstlisting}[language=C++]
#include <pluss/pluss.h>
// narrower:
#include <pluss/coden/coden.h>
#include <pluss/coden/note_factory.h>

auto s = pluss::content("hello");
auto v = pluss::content({"a", "b", "c"});
auto m = pluss::content({{"a","b"}, {"c","d"}});
\end{lstlisting}

\subsection{Factory Patterns You Will Actually Use}
\begin{lstlisting}[language=C++]
std::vector<std::string> flat{"p","q","r","s"};
auto n = pluss::make_note(flat, std::array<std::int64_t,2>{2,2});

auto from_lines = pluss::note_from_lines("line1\nline2\nline3");
auto from_words = pluss::note_from_words("alpha beta gamma");

auto filled = pluss::full_note({2,1,2}, "z");
auto empty  = pluss::empty_note({2,2});

// raw bytes mode
auto bytes = pluss::bytes_content({{"\x01","\x02"}, {"\x03","\x04"}});
\end{lstlisting}

\subsection{Common Workflows}
\subsubsection{Text cleanup + normalization}
\begin{lstlisting}[language=C++]
auto docs = pluss::content("One.\nTwo!");
auto clean = docs.lower().collapse_whitespace().strip();
\end{lstlisting}

\subsubsection{Tokenization and split}
\begin{lstlisting}[language=C++]
auto docs = pluss::content("alpha beta gamma");
auto words = docs.tokenize_words();
for (std::int64_t i = 0; i < words.tokens().numel(); ++i) {
  std::cout << words.tokens().at(i) << "\n";
}
\end{lstlisting}

\subsubsection{NLP/TeX utilities}
\begin{lstlisting}[language=C++]
auto n = pluss::content("Email bob@example.com at 9:30am on 2025-01-03");
auto masked = n.nlp_mask_pii();

auto tex = pluss::content("\\textbf{Bold} and $x+y$");
auto txt = tex.tex_to_text();
\end{lstlisting}

\subsubsection{Printing control}
\begin{lstlisting}[language=C++]
pluss::NotePrintOptions opt;
opt.max_cell_chars = 24;
opt.max_elements = 10;
opt.edge_items = 2;
std::cout << pluss::print_note(docs, opt) << "\n";
\end{lstlisting}

\subsection{Note Demo Map}
\begin{longtable}{@{}p{0.36\linewidth}p{0.56\linewidth}@{}}
\textbf{Demo} & \textbf{Best for} \\
\hline
\texttt{note\_factory\_demo.cpp} & All construction/factory paths \\
\texttt{note\_base\_methods\_demo.cpp} & Shape/index/iterator/base helpers \\
\texttt{note\_method\_families\_demo.cpp} & NLP/TeX/token method families \\
\texttt{note\_printing\_demo.cpp} & Display and bytes rendering options \\
\end{longtable}

\section{Timeline and Duration (CodeH)}
\subsection{Construct Timeline}
\begin{lstlisting}[language=C++]
#include <pluss/pluss.h>
// narrower:
#include <pluss/codeh/codeh.h>
#include <pluss/codeh/timeline_factory.h>

auto t0 = pluss::moment("2024-01-01T00:00:00Z");
auto t1 = pluss::moment({"2024-01-01T00:00:00Z", "2024-01-01T01:00:00Z"});
auto t2 = pluss::moment({{"2024-02-01T00:00:00Z", "2024-02-01T01:00:00Z"}});

auto t3 = pluss::from_epoch_nanos(std::vector<std::int64_t>{0, 1'000'000'000}, std::vector<std::int64_t>{2});
\end{lstlisting}

\subsection{Construct Duration and do date arithmetic}
\begin{lstlisting}[language=C++]
auto d1 = pluss::duration_seconds(30);
auto d2 = pluss::duration_minutes(5);
auto d3 = pluss::duration_hours(1);
auto d4 = pluss::duration_days(1);

auto t = pluss::moment("2024-01-01T00:00:00Z");
auto t_plus = t + d2;
auto diff = t_plus - t;   // Duration
\end{lstlisting}

\subsection{Datetime field extraction, formatting, timezone}
\begin{lstlisting}[language=C++]
auto t = pluss::moment({"2024-01-01T00:00:00Z", "2024-06-15T12:30:00Z"});
auto years = t.year();
auto months = t.month();
auto weekdays = t.weekday();

auto text = t.strftime("%Y-%m-%d %H:%M:%S");
auto local = t.tz_convert("UTC");
auto floored = t.floor("D");
\end{lstlisting}

\subsection{Timeline Demo Map}
\begin{longtable}{@{}p{0.36\linewidth}p{0.56\linewidth}@{}}
\textbf{Demo} & \textbf{Best for} \\
\hline
\texttt{timeline\_non\_generated\_and\_factory\_methods\_demo.cpp} & Factories + legacy/non-generated helpers + conversions \\
\texttt{timeline\_printing\_2d\_3d\_demo.cpp} & 2D/3D display and print options \\
\end{longtable}

\section{TimeIndex (CodeH)}
\subsection{Construct and inspect}
\begin{lstlisting}[language=C++]
#include <pluss/pluss.h>
#include <pluss/codeh/timeindex_factory.h>

auto idx = pluss::make_timeindex({
  "2024-01-01T00:00:00Z",
  "2024-01-01T00:15:00Z",
  "2024-01-01T00:45:00Z",
  "2024-01-01T01:05:00Z",
});

std::cout << idx.freq() << "\n";
std::cout << idx.is_monotonic_increasing() << "\n";
\end{lstlisting}

\subsection{Resample and rolling}
\begin{lstlisting}[language=C++]
auto values = pluss::value(std::vector<double>{10,20,30,40}, {4});

auto hourly_sum = idx.resample_agg(values, "H", "sum", 1);
auto hourly_mean = idx.resample_agg(values, "H", "mean", 1);

auto win = pluss::duration_minutes(45);
auto roll_sum = idx.rolling_agg(values, win, "sum");
auto roll_max = idx.rolling_agg(values, win, "max");
\end{lstlisting}

\subsection{Slicing, asof, set-like ops}
\begin{lstlisting}[language=C++]
auto k = pluss::moment("2024-01-01T00:40:00Z");
auto p = idx.searchsorted(k, false);
auto a = idx.asof(k);
auto s = idx.slice_time("2024-01-01T00:00:00Z", "2024-01-01T01:00:00Z");
\end{lstlisting}

\subsection{TimeIndex Demo Map}
\begin{longtable}{@{}p{0.36\linewidth}p{0.56\linewidth}@{}}
\textbf{Demo} & \textbf{Best for} \\
\hline
\texttt{timeline\_dataframe\_timeindex\_methods\_demo.cpp} & Frequency inference, resample, rolling aggregation \\
\end{longtable}

\section{Method Families (Quick Reference)}
\subsection{Series}
\begin{itemize}[leftmargin=1.5em]
  \item Core tensor ops: shape/view/permute/select/slice
  \item Numeric ops: arithmetic, unary math, comparison, logical, bitwise
  \item Reductions: sum/mean/min/max/var/std/norm/all/any
  \item Data movement: to(platform), to(dtype), contiguous, clone, to\_host
  \item Indexing/masking: where, take, gather/scatter/scatter\_add/scatter\_reduce, masked\_select
  \item Stats/signal extras: pdf/cdf/histcounts/skewness/kurtosis/dct/dst/hilbert
\end{itemize}

\subsection{Note}
\begin{itemize}[leftmargin=1.5em]
  \item Text normalization/cleanup, regex, replace/remove families
  \item Tokenization families: split/tokenize/ngrams/extractors
  \item NLP families: masking, stopwords, language/readability/sentiment helpers
  \item TeX families: strip/normalize/extract/convert helpers
\end{itemize}

\subsection{Timeline/TimeIndex}
\begin{itemize}[leftmargin=1.5em]
  \item Timeline: datetime field extraction, predicates, tz, rounding, format, min/max/sort/unique
  \item TimeIndex: monotonic/unique checks, search/slice/asof, resample, rolling, set operations
\end{itemize}

\section{Exhaustive API Index Files}
\label{sec:api-index-files}
These files are generated from code and are intended as searchable reference inventories:
\begin{itemize}[leftmargin=1.5em]
  \item \texttt{docs/generated/series\_methods\_generated.txt}
  \item \texttt{docs/generated/note\_methods\_decls\_from\_code.h}
  \item \texttt{docs/generated/timeline\_methods\_signatures\_from\_code.h}
  \item \texttt{docs/generated/timeindex\_methods\_signatures\_from\_code.h}
\end{itemize}

\section{Sphinx Migration Note}
For long-term user-facing docs, Sphinx is a better fit than a single LaTeX file:
\begin{itemize}[leftmargin=1.5em]
  \item Multi-page structure with navigation/search
  \item Better API reference layering (tutorials + reference + examples)
  \item Easier incremental updates as methods evolve
\end{itemize}
A starter Sphinx skeleton is included under \texttt{docs/sphinx/}.

\section{Build This PDF}
\begin{lstlisting}[language=bash]
pdflatex -output-directory docs docs/PLUSSMA_USER_MANUAL.tex
pdflatex -output-directory docs docs/PLUSSMA_USER_MANUAL.tex
\end{lstlisting}

\end{document}
